\subsection{Lists and Counting Principles}

Counting translates structured finite choices into cardinalities.
In uniform finite probability models, those cardinalities later determine probabilities.

\begin{definition} \label{def:finite-sets-and-standard-index-sets} \hlt{Finite Sets and Standard Index Sets}\\
Set $J_0\defeq\emptyset$, and for $n\in\N$ set $J_n\defeq\{1,\ldots,n\}$.
A set $A$ is \hlt{finite} if it admits a bijection $e:J_n\to A$ with $n\in\N_0$.
\end{definition}

\begin{proposition} \label{prop:finite-cardinality-well-defined} \hlt{Finite Cardinality Is Well Defined}\\
Let $A$ be a set, let $m,n\in\N_0$, and let $e:J_m\to A$ and $g:J_n\to A$ be bijections.
Then $m=n$.
Consequently, if $A$ is finite, define $\#A$ as the unique $n\in\N_0$ for which a bijection $J_n\to A$ exists.
\begin{proof}
It suffices to prove that a bijection $J_m\to J_n$ forces $m=n$, because $g^{-1}\circ e:J_m\to J_n$ is bijective.
For $r\in\N$, let $P_r$ assert that a bijection $J_{r-1}\to J_n$, with $n\in\N_0$, can exist only when $r-1=n$.
We apply Theorem \ref{thm:mathematical-induction}.
For $r=1$, the domain $J_0$ is empty, so a bijection $J_0\to J_n$ forces $n=0$.
Assume $P_r$ and let $\varphi:J_r\to J_n$ be bijective.
Since $J_r\neq\emptyset$, $n\in\N$.
Put $q\defeq\varphi(r)$.
The restriction $\varphi|_{J_{r-1}}$ is a bijection $J_{r-1}\to J_n\setminus\{q\}$.
Define $\rho:J_{n-1}\to J_n\setminus\{q\}$ by $\rho(i)=i$ for $i<q$ and $\rho(i)=i+1$ for $i\geq q$.
Then $\rho$ is bijective, so $\rho^{-1}\circ\varphi|_{J_{r-1}}:J_{r-1}\to J_{n-1}$ is bijective.
By $P_r$, $r-1=n-1$, and hence $r=n$.
Thus $P_{r+1}$ holds, and $(\forall r\in\N)\,P_r$ by induction.\\
\hlt{Structural Commentary}: Induction removes one matched terminal point from two finite index sets.\\
\hlt{Key Insight}: A bijection preserves the number of listed positions.\\
\end{proof}
\end{proposition}

\begin{definition} \label{def:lists-words-and-product-sets} \hlt{Lists, Words, and Product Sets}\\
Let $p\in\N$ and let $S_1,\ldots,S_p$ be sets.
\begin{enumerate}[label=(\roman*)]
\setlength{\itemsep}{0pt}
\item A \hlt{list} of length $p$ with $i$th entry in $S_i$ is an element $(s_1,\ldots,s_p)\in S_1\times\cdots\times S_p$.
\item If $S_1=\cdots=S_p=S$, a list in $S^p$ is a \hlt{word} of length $p$ over the alphabet $S$.
\item A counting model is \hlt{ordered} when its outcomes are lists and coordinate order is retained.
Thus interchanging two unequal entries produces a different outcome.
\end{enumerate}
\end{definition}

\begin{proposition} \label{prop:addition-subtraction-and-division-principles} \hlt{Addition, Subtraction, and Division Principles}\\
\begin{enumerate}[label=(\roman*)]
\setlength{\itemsep}{0pt}
\item \hlt{Addition}: Let $r\in\N$ and let $A_1,\ldots,A_r$ be pairwise disjoint finite sets.
Then $\bigcup\limits_{i=1}^{r}A_i$ is finite and $\#\bigl(\bigcup\limits_{i=1}^{r}A_i\bigr)=\sum\limits_{i=1}^{r}\#A_i$.
\item \hlt{Subtraction}: If $A$ is finite and $B\subseteq A$, then $B$ and $A\setminus B$ are finite and $\#(A\setminus B)=\#A-\#B$.
\item \hlt{Division}: Let $f:Y\to X$ be a surjection between finite sets.
Suppose its fibers are finite and $d\in\N$ satisfies $(\forall x\in X)\,\#f^{-1}(\{x\})=d$.
Then $\#Y=d\#X$ and $\#X=\frac{\#Y}{d}$.
\end{enumerate}
\begin{proof}
\begin{enumerate}[label=(\roman*)]
\setlength{\itemsep}{0pt}
\item Put $n_i\defeq\#A_i$, choose a bijection $e_i:J_{n_i}\to A_i$, and set $q_0\defeq0$ and $q_k\defeq\sum\limits_{i=1}^{k}n_i$ for $k\in J_r$.
For $i\in J_r$, set $I_i\defeq\{t\in J_{q_r}:q_{i-1}<t\leq q_i\}$.
The sets $I_i$ partition $J_{q_r}$.
For $t\in I_i$, set $e(t)\defeq e_i(t-q_{i-1})$.
This defines a map $e:J_{q_r}\to\bigcup\limits_{i=1}^{r}A_i$.
The disjointness of the $A_i$ and the bijectivity of each $e_i$ make $e$ bijective.
Thus the union is finite and has cardinality $q_r=\sum\limits_{i=1}^{r}\#A_i$.
\item For $r\in\N$, let $P_r$ assert the claim whenever $\#A=r-1$.
We apply Theorem \ref{thm:mathematical-induction}.
If $r=1$, then $A=\emptyset$, so the claim holds.
Assume $P_r$, let $\#A=r$, choose a bijection $e:J_r\to A$, put $a\defeq e(r)$, and set $A'\defeq A\setminus\{a\}$.
Then $e|_{J_{r-1}}:J_{r-1}\to A'$ is bijective, so $\#A'=r-1$.
If $a\in B$, apply $P_r$ to $B\setminus\{a\}\subseteq A'$ and use the disjoint union $B=(B\setminus\{a\})\cup\{a\}$.
If $a\notin B$, apply $P_r$ to $B\subseteq A'$ and use the disjoint union $A\setminus B=(A'\setminus B)\cup\{a\}$.
In either case, part (i) proves that both sets are finite and that $\#A=\#B+\#(A\setminus B)$.
Rearranging proves the asserted formula and completes the induction.
\item Put $m\defeq\#X$.
If $m=0$, surjectivity forces $Y=\emptyset$, and both conclusions follow.
If $m\in\N$, choose a bijection $e:J_m\to X$ and write $x_i\defeq e(i)$.
The fibers $f^{-1}(\{x_i\})$ are pairwise disjoint and have union $Y$.
By part (i) and Proposition \ref{prop:finite-sum-basic-laws}, $\#Y=\sum\limits_{i=1}^{m}d=d\sum\limits_{i=1}^{m}1=dm=d\#X$.
Since $d\neq0$, division gives $\#X=\frac{\#Y}{d}$.
\end{enumerate}
\hlt{Structural Commentary}: Explicit listings reduce unions, complements, and fibers to disjoint blocks.\\
\hlt{Key Insight}: Cardinalities add over disjoint representations.\\
\end{proof}
\end{proposition}

\begin{theorem} \label{thm:multiplication-principle-finite-products} \hlt{Multiplication Principle}\\
Let $p\in\N$ and let $n_1,\ldots,n_p\in\N_0$.
\begin{enumerate}[label=(\roman*)]
\setlength{\itemsep}{0pt}
\item \hlt{Finite products}: If $S_1,\ldots,S_p$ are finite and $(\forall i\in J_p)\,\#S_i=n_i$, then $\#(S_1\times\cdots\times S_p)=\prod\limits_{i=1}^{p}n_i$.
\item \hlt{Successive choices}: Set $H_0\defeq\{\emptyset\}$.
Recursively, for $i\in J_p$ and $h\in H_{i-1}$, let $C_i(h)$ be finite with $\#C_i(h)=n_i$, and set $H_i\defeq\{(h,c):h\in H_{i-1},\ c\in C_i(h)\}$.
Then $H_p$ is finite and $\#H_p=\prod\limits_{i=1}^{p}n_i$.
\end{enumerate}
\begin{proof}
First let $A,B$ be finite.
If $A=\emptyset$, then $A\times B=\emptyset$.
Otherwise put $m\defeq\#A$, choose a bijection $a:J_m\to A$, and write $a_i\defeq a(i)$.
The sets $\{a_i\}\times B$ are pairwise disjoint, each is bijective with $B$, and their union is $A\times B$.
By Proposition \ref{prop:addition-subtraction-and-division-principles}, $\#(A\times B)=m\#B=\#A\#B$.
\begin{enumerate}[label=(\roman*)]
\setlength{\itemsep}{0pt}
\item Apply Theorem \ref{thm:mathematical-induction} to $p$.
For $p=1$, the map $s\mapsto(s)$ is a bijection $S_1\to S_1^1$, so the formula holds.
Assume the formula for $p$ factors.
Flattening the final ordered pair identifies $(S_1\times\cdots\times S_p)\times S_{p+1}$ bijectively with $S_1\times\cdots\times S_{p+1}$.
The binary identity and the induction hypothesis give $\#(S_1\times\cdots\times S_{p+1})=(\prod\limits_{i=1}^{p}n_i)n_{p+1}=\prod\limits_{i=1}^{p+1}n_i$.
\item The set $H_0$ is finite with $\#H_0=1$.
Suppose $H_{i-1}$ is finite.
If $H_{i-1}=\emptyset$, then $H_i=\emptyset$.
Otherwise the sets $\{h\}\times C_i(h)$, indexed by $h\in H_{i-1}$, are pairwise disjoint, have cardinality $n_i$, and have union $H_i$.
By Proposition \ref{prop:addition-subtraction-and-division-principles}, $\#H_i=n_i\#H_{i-1}$ in both cases.
Applying Theorem \ref{thm:mathematical-induction} to these base and step clauses gives $\#H_p=n_p\cdots n_1=\prod\limits_{i=1}^{p}n_i$.
\end{enumerate}
\hlt{Structural Commentary}: Disjoint extension fibers convert each additional choice stage into multiplication.\\
\hlt{Key Insight}: Uniform branching multiplies the number of histories.\\
\end{proof}
\end{theorem}

\begin{example} \label{ex:cards-binary-and-dna-counting} \hlt{Finite Words by the Multiplication Principle}\\
By Theorem \ref{thm:multiplication-principle-finite-products}, the following slot counts multiply.
\begin{enumerate}[label=(\roman*)]
\setlength{\itemsep}{0pt}
\item A playing card is determined by one of $4$ suits and one of $13$ ranks, giving $4\cdot13=52$ cards.
\item The binary alphabet $\{0,1\}$ has $\#(\{0,1\}^8)=2^8=256$ words of length $8$.
\item For $N\in\N$, the alphabet $\{A,C,G,T\}$ has $4^N$ words of length $N$; in particular, there are $4^3=64$ codons.
\end{enumerate}
\hlt{Key Insight}: Independent slots multiply; word counts grow exponentially with length.\\
\end{example}
